\subsection{External benchmarks}
\label{sec:exp_external}
Both external evaluations below were run with the earlier box-policy front end (\cref{tab:components}); re-running them with the final pipeline is pending, and we report them as they are.

\paragraph{Editing correct text (VTE-Bench).} On the 100 real clips of \steer's own benchmark~\cite{zeng2026steervte}, where correct text must be changed into a different string, \steer reaches 92\% readability and our earlier pipeline 82\% (8 vs.\ 18 discordant clips, $p=0.076$; a second reader, EasyOCR, gives 49\% vs.\ 44\%). Editing differs from repair: the source text is valid and must be removed entirely, which is the regime that generative editors are trained for and in which our gates provide no advantage. We include the result to delimit our claims.

\paragraph{Commercial video models (VidScribe).} VidScribe~\cite{vidscribe2026} releases the outputs of seven commercial video models for video-to-video text editing. On the 44 of its 93 clips that have a localisable text region, our earlier pipeline reaches 29.5\% (naive paste) and 27.3\% (box policy) readability, ahead of the strongest commercial model (minimax-h3, 18.2\%) and wan-3.0 (15.9\%), although not significantly (box policy vs.\ minimax-h3: 6 vs.\ 2 discordant clips, $p=0.29$); against the other five models (0--13.6\%) the box policy is significantly better ($p\le0.031$). Commercial models regenerate the whole clip and lower the PSNR outside the text region to 19.8--27.9\,dB, against 37.6--37.9\,dB for our outputs. Absolute readability is low for every method because most VidScribe targets are long phrases.
